\indent ಸಿದ್ಧಿಸುವುವು ; ಇವು ಲಿಂಗ ವಚನ, ವಿಭಕ್ತಿಗಳಿಂದ ವ್ಯತ್ಯಾಸವಾಗದ ರೂಪಗಳು.\par

\vspace{3pt plus 2pt}\noindent\hspace*{1.2em}\kbvlabel{ಉದಾ}\par\vspace{1pt}

\begin{kbvderiv}
\vspace{4pt plus 3pt minus 1pt}
\begin{center}\begin{adjustbox}{max width=\textwidth}
{\renewcommand{\arraystretch}{1.22}\setlength{\tabcolsep}{4pt}
\begin{tabular}{l c l}
೧. ಬೆಳಕು ಬರಲು ಕತ್ತಲೆ ಸರಿಯಿತು. (ಬರಲು & {\color{kbvgrey}$=$} & ಬರು + ಅಲು) \\
೨. ಈ ಹಣ್ಣು ತಿನ್ನಲಿಕ್ಕೆ ರುಚಿಯಾಗಿದೆ. (ತಿನ್ನಲಿಕ್ಕೆ & {\color{kbvgrey}$=$} & ತಿನ್ನು + ಅಲಿಕ್ಕೆ) \\
೩. ನೀನು ರೊಟ್ಟಿ ತಿಂದು ಹೋಗು. (ತಿಂದು & {\color{kbvgrey}$=$} & ತಿನ್ನು + ದು) \\
೪. ಮಾಡಿ ನೋಡು ಮದ್ದಿನ ಗುಣ. (ಮಾಡಿ & {\color{kbvgrey}$=$} & ಮಾಡು + ಇ) \\
೫. ಊರಿಗೆ ಹೋದರೆ ಹಿಂದಿರುಗಲಾರೆ. (ಹೋದರೆ & {\color{kbvgrey}$=$} & ಹೋಗು + ದರೆ) \\
೬. ಅವಳು ನಗುತ್ತೆ ಮನೆಗೆ ಹೋದಳು. (ನಗುತ್ತೆ & {\color{kbvgrey}$=$} & ನಗು + ಉತ್ತೆ) \\
೭. ನಾನು ನಿನ್ನನ್ನು ನೋಡದೆ ಇರಲಾರೆ. (ನೋಡದೆ & {\color{kbvgrey}$=$} & ನೋಡು + ಅದೆ) \\
೮. ಕಮಲೆಯು ನನ್ನನ್ನು ಬರೆ ಹೇಳಿದ್ದಾಳೆ. (ಬರೆ & {\color{kbvgrey}$=$} & ಬರು + ಅ) \\
೯. ಕಿಂದರಿ ಜೋಗಿ ಹಾಡುತೆ ಬರುತಿಹನು. (ಹಾಡುತೆ & {\color{kbvgrey}$=$} & ಹಾಡು + ಉತ್ತೆ) \\
\end{tabular}}
\end{adjustbox}\end{center}
\vspace{4pt plus 3pt minus 1pt}
\end{kbvderiv}

\vspace{2.5pt plus 3pt minus 0.5pt}\noindent\hspace*{1.2em}\kbvnum{೧೦.}\hspace{0.55em}ರಾಷ್ಟ್ರಾಧ್ಯಕ್ಷರು ಬರುತ್ತಿರಲು ಮಕ್ಕಳು ಚಪ್ಪಾಳೆ ತಟ್ಟಿದರು.\par

\indent (ಬರುತ್ತಿರಲು = ಬರು + ಉತ್ತೆ + ಅರಲು)\par

\indent ನೆನಪಿಡಿ : ಕೃದಂತಾವ್ಯಯಗಳು ಪೂರ್ಣ ಕ್ರಿಯಾರ್ಥವನ್ನು ಕೊಡುವುದಿಲ್ಲವಾದ ಕಾರಣ ಮತ್ತೊಂದು ಕ್ರಿಯೆಯನ್ನು ಅಪೇಕ್ಷಿಸುವುದು. ಇವುಗಳ ಲಕ್ಷಣ, ಉದಾ: ನಗುತ್ತ ಇದ್ದನು; ನೋಡದೆ ಇರಲಾರೆ; ಮಾಡಿ ನೋಡು.\par

\kbvexercise{ಅಭ್ಯಾಸ ಪತ್ರಿಕೆ ೧೪}

\vspace{2.5pt plus 3pt minus 0.5pt}\noindent\hspace*{1.2em}\kbvnum{1.}\hspace{0.55em}“ಬರೆಯುವನು” ಈ ಕ್ರಿಯಾಪದದಲ್ಲಿ ಸೇರಿರುವ ಕಾಲಸೂಚಕ ಪ್ರತ್ಯಯ ಯಾವುದು? ಮತ್ತು ಆಖ್ಯಾತ ಪ್ರತ್ಯಯ ಯಾವುದು?\par

\vspace{2.5pt plus 3pt minus 0.5pt}\noindent\hspace*{1.2em}\kbvnum{2.}\hspace{0.55em}ಸ್ತ್ರೀಲಿಂಗ ಏಕವಚನದಲ್ಲಿ “ಕೊಡು” ಧಾತುವಿನ ಭವಿಷ್ಯಕಾಲದ ಕ್ರಿಯಾಪದ ರೂಪವನ್ನು ಬರೆಯಿರಿ.\par

\vspace{2.5pt plus 3pt minus 0.5pt}\noindent\hspace*{1.2em}\kbvnum{3.}\hspace{0.55em}“ಹೋದನು” ಈ ಕ್ರಿಯಾಪದವನ್ನು ಬಿಡಿಸಿ ಬರೆಯಿರಿ.\par

\vspace{2.5pt plus 3pt minus 0.5pt}\noindent\hspace*{1.2em}\kbvnum{4.}\hspace{0.55em}“ಕಾಣುವ” ಈ ಶಬ್ದದಲ್ಲಿ ಸೇರಿರುವ ಪ್ರತ್ಯಯಗಳು ಯಾವುವು? ಅವುಗಳ ವಿಶೇಷವೇನು?\par

\vspace{2.5pt plus 3pt minus 0.5pt}\noindent\hspace*{1.2em}\kbvnum{5.}\hspace{0.55em}ಕೃದಂತ ನಾಮ ಎಂದರೇನು? ಉದಾಹರಣೆಗಳ ಸಹಿತ ತಿಳಿಸಿ.\par
